\documentclass{article}
\usepackage{amsmath}
\newtheorem{assumption}{Assumption}

\title{A Worked Gaussian Integral}
\author{A. Nonymous}

\begin{document}
\maketitle

% A comment that must not reach the extracted prose.
\section{Setup}

We consider the one-dimensional Gaussian integral, which appears whenever a quadratic
action is integrated out.

\begin{assumption}\label{as:positive}
The parameter $a$ is real and strictly positive, so the integral converges.
\end{assumption}

Under Assumption~\ref{as:positive} the integral is well defined and we write

\begin{equation}\label{eq:definition}\tag{1}
I(a) = \int_{-\infty}^{\infty} e^{-a x^{2}}\,dx
\end{equation}

where $a$ is the inverse width of the Gaussian, and $x$ is the integration variable.

\section{Evaluation}

Squaring \eqref{eq:definition} and renaming the variable in the second factor gives

\begin{equation}\label{eq:squared}\tag{2}
I(a)^{2} = \int_{-\infty}^{\infty}\int_{-\infty}^{\infty} e^{-a(x^{2}+y^{2})}\,dx\,dy
\end{equation}

which is a two-dimensional integral over the whole plane.

Changing to polar coordinates, so that $r$ is the radial coordinate and $\theta$ is the
polar angle, we obtain

\begin{equation}\label{eq:polar}\tag{3}
I(a)^{2} = \int_{0}^{2\pi}\int_{0}^{\infty} e^{-a r^{2}} r\,dr\,d\theta
\end{equation}

The angular integral contributes a factor of $2\pi$ and the radial integral is elementary,
so that

\begin{equation}\label{eq:result}\tag{4}
I(a) = \sqrt{\frac{\pi}{a}}
\end{equation}

which is the standard result quoted in \eqref{eq:definition}.

An unnumbered aside, which the extractor should keep but not number:

\begin{equation*}
\lim_{a \to \infty} I(a) = 0
\end{equation*}

\end{document}
